% =====================================================================
% Act I --- Compression
% Scenes 4 through 12 of the W2 narrative.
% Total: ~12 slides + 1 optional act-opener divider.
% =====================================================================

% Set the section indicator for every slide in this file.
\renewcommand{\AINITsection}{Act I \textemdash{} Compression}

% ---------------------------------------------------------------------
% Act-opener divider (optional but pedagogically nice).
% ---------------------------------------------------------------------
\sectiondivider%
  {``\textellipsis{} a Map of the Empire whose size was that of the Empire, and which coincided point for point with it.''}%
  {Jorge Luis Borges, \textit{On Exactitude in Science}, 1946}%
  {Act I \textemdash{} Compression}

% ---------------------------------------------------------------------
% Scene 4a --- The data
% ---------------------------------------------------------------------
\begin{frame}{The data: 21 readers, 16 books}
  A small public library has kept a quiet record.

  \vspace{0.5em}

  For each of \AINITemph{21 anonymous readers} and each of \AINITemph{16 books},
  one cell records whether the reader liked the book, disliked it, or
  has not yet read it.

  \vspace{0.5em}

  Three symbols. Three hundred and fifty-seven cells.

  \vspace{0.5em}

  No names, no genres, no author photographs --- only the matrix.
\end{frame}

% ---------------------------------------------------------------------
% Scene 4b --- The matrix as raw data, with a question
% ---------------------------------------------------------------------
\begin{frame}{Raw data}
  \begin{center}
    \begin{tikzpicture}[every node/.style={inner sep=0.6pt}]
      % miniature 21x16 grid of dots, three shades, schematic.
      \foreach \r in {0,...,20} {
        \foreach \c in {0,...,15} {
          \pgfmathtruncatemacro{\h}{mod(\r*7+\c*3,11)}
          \ifnum\h<4
            \fill[ink!85] (\c*0.26,-\r*0.18) circle (1.3pt);
          \else
            \ifnum\h<7
              \fill[mutedWarmGray] (\c*0.26,-\r*0.18) circle (0.8pt);
            \else
              \draw[mutedWarmGray!50] (\c*0.26,-\r*0.18) circle (1.1pt);
            \fi
          \fi
        }
      }
      \node[font=\itshape\small,mutedWarmGray] at (2.0,0.5) {16 books $\rightarrow$};
      \node[font=\itshape\small,mutedWarmGray,rotate=90] at (-0.7,-1.8) {$\leftarrow$ 21 readers};
    \end{tikzpicture}
  \end{center}

  \begin{center}
    \itshape Three hundred and fifty-seven cells. Can your brain remember them?
  \end{center}
\end{frame}

% ---------------------------------------------------------------------
% Scene 5 --- Block-diagonal puzzle (interactive)
% ---------------------------------------------------------------------
\vizhandoff{01\_block-diagonal-puzzle.html}

% ---------------------------------------------------------------------
% Scene 6a --- The four-corner reveal: the axes named
% ---------------------------------------------------------------------
\begin{frame}{The geometry was always there}
  \sidenote{Both axes are old library science: Dewey decimal (1876) gives Fiction vs.\ Nonfiction; Flesch (1948) and Lexile (1989) give Light vs.\ Demanding.}
  When the matrix is sorted by hand, two axes appear:

  \vspace{0.6em}

  \begin{itemize}
    \setlength{\itemsep}{0.4em}
    \item \AINITemph{Fiction} versus \AINITemph{Nonfiction}
    \item \AINITemph{Light} versus \AINITemph{Demanding}
  \end{itemize}

  \vspace{0.8em}

  Neither axis was invented for this dataset. They are how librarians and
  educators have sorted books for a century.

  \vspace{0.4em}

  We did not impose them. We \AINITemph{recovered} them.
\end{frame}

% ---------------------------------------------------------------------
% Scene 6b --- The 16 books placed in 4 quadrants
% ---------------------------------------------------------------------
\begin{frame}{Sixteen books, four quadrants}
  \begin{center}
    \begin{tikzpicture}[
      every node/.style={font=\footnotesize},
      book/.style={font=\itshape\footnotesize,inner sep=1pt},
    ]
      % axes
      \draw[ink,thick,->] (-4.2,0) -- (4.2,0) node[right,font=\scshape\scriptsize,mutedWarmGray] {nonfiction};
      \draw[ink,thick,->] (0,-2.8) -- (0,2.8) node[above,font=\scshape\scriptsize,mutedWarmGray] {demanding};
      \node[left,font=\scshape\scriptsize,mutedWarmGray] at (-4.2,0) {fiction};
      \node[below,font=\scshape\scriptsize,mutedWarmGray] at (0,-2.8) {light};

      % quadrant tints
      \fill[accentOxblood,opacity=0.04] (-4,0) rectangle (0,2.6);
      \fill[accentOxblood,opacity=0.04] (0,-2.6) rectangle (4,0);

      % Books and quadrants per BOOKS_MATRIX.md (the matrix that drives every viz).
      % top-left: demanding fiction (DF)
      \node[book] at (-3.0, 2.1) {Ulysses};
      \node[book] at (-1.5, 1.6) {Borges};
      \node[book] at (-2.7, 1.1) {War \& Peace};
      \node[book] at (-1.2, 0.5) {Lighthouse};

      % top-right: demanding nonfiction (DN)
      \node[book] at ( 1.2, 2.3) {Russell};
      \node[book] at ( 2.7, 1.7) {Kant};
      \node[book] at ( 1.6, 1.0) {Smith};
      \node[book] at ( 3.2, 0.5) {Darwin};

      % bottom-left: light fiction (LF)
      \node[book] at (-2.8,-0.6) {Christie};
      \node[book] at (-1.4,-1.2) {Potter};
      \node[book] at (-3.1,-1.8) {Hunger Games};
      \node[book] at (-1.2,-2.3) {Da Vinci Code};

      % bottom-right: light nonfiction (LN)
      \node[book] at ( 1.4,-0.6) {Eat Pray Love};
      \node[book] at ( 3.0,-1.1) {Joy of Cooking};
      \node[book] at ( 1.7,-1.9) {7 Habits};
      \node[book] at ( 3.1,-2.3) {Switzerland};
    \end{tikzpicture}
  \end{center}
\end{frame}

% ---------------------------------------------------------------------
% Scene 7 --- Latent axes lab (interactive)
% ---------------------------------------------------------------------
\vizhandoff{02\_latent-axes-lab.html}

% ---------------------------------------------------------------------
% Scene 8 --- Dot product equation
% ---------------------------------------------------------------------
\equationslide{Preference is a dot product}%
  {\hat r_{ui} = u_i^\top v_j}%
  {The user's vector and the book's vector live in the same space; their inner product is the predicted preference.}

% ---------------------------------------------------------------------
% Scene 9 --- Spectral microscope (interactive)
% ---------------------------------------------------------------------
\vizhandoff{spectral-microscope.html}

% ---------------------------------------------------------------------
% Scene 10 --- Eckart & Young (1936)
% ---------------------------------------------------------------------
\begin{frame}{Eckart \& Young (1936)\textsuperscript{1}}
  \begin{quote}
    \itshape Among all rank-$k$ approximations of any matrix, the one
    given by the top-$k$ singular vectors is optimal in squared error.
  \end{quote}

  \vspace{0.8em}

  No algorithm beats it. No clever trick beats it. No future
  architecture beats it.

  \vspace{0.4em}

  The curve you just saw is not an empirical observation. It is the
  graph of a theorem ninety years old.
\end{frame}

% ---------------------------------------------------------------------
% Scene 11 --- Customer-Urns Carousel (interactive)
% ---------------------------------------------------------------------
\vizhandoff{customer-urns-carousel.html}

% ---------------------------------------------------------------------
% Scene 12 --- Embeddings that learn (interactive)
% ---------------------------------------------------------------------
\vizhandoff{07\_embeddings-that-learn.html}

% ---------------------------------------------------------------------
% Closing slide for Act I
% ---------------------------------------------------------------------
\begin{frame}{End of Act I}
  We compressed a matrix. We named the act.

  \vspace{0.8em}

  The names --- \AINITemph{entropy}, \AINITemph{rank},
  \AINITemph{factorization} --- were old words from 1936 and 1948.
\end{frame}
